% Appendix A: Floating-point essentials (fundamentals author). PLAN.md section 1.4.
% Uses T16 (tab:fp-constants), typeset from the DS17 values of data/values.tex.
% Language pass for students (PLAN.md section 12.9): the wording only; content, evidence and order as at 3192e8a.
\chapter{Floating-point essentials}\label{ch:floating-point}

This appendix collects the facts about floating-point arithmetic that \cref{ch:rounding} and the sidebars use. We
start with how a number is stored. Then we look at the spacing of the numbers and at rounding, at the two ends of the
range, and at what rounding does to $x + h$ and to a difference. The appendix ends with other number types. The
standard references are Higham's book and Goldberg's survey.\cite{higham2002,goldberg1991} The formats and the
rounding rule are fixed by the IEEE~754 standard.\cite{ieee754-2019}

\section{Representation}\label{sec:fp-representation}

A binary floating-point number is
\begin{equation*}
  x = \pm m \cdot 2^{e}, \qquad 1 \le m < 2,
\end{equation*}
where the \emph{significand} $m$ has $p$ binary digits, and the \emph{exponent} $e$ is an integer in a fixed range.
This is scientific notation in base 2: the significand holds the digits, and the exponent says where the binary point
goes. For \code{float}, $p = \num{\valDigitsFloat}$; for \code{double}, $p = \num{\valDigitsDouble}$
(\cref{tab:fp-constants}).

Every integer up to $2^{p}$ is exact, that is, the type stores it without rounding. This is why the library's integer
weights are exact in every type with at least 24 significant bits: they are all below $2^{24}$. Between two powers of
two, the numbers are evenly spaced, because the significand always has the same $p$ digits and the exponent is
fixed. At each power of two, the exponent goes up by one, and the spacing doubles.

\begin{table}
\sidecaption{The constants of \code{float} and \code{double}, read from \code{std::numeric_limits} on the build
machine (IEEE~754 formats). Here $\epsmach = 2^{1-p}$. The decimal digits are \code{digits10}, the number of decimal
digits that survive a round trip, from decimal text to the type and back. Below the smallest normal number, the
subnormal numbers continue down to the smallest subnormal.\label{tab:fp-constants}}
\centering\small
\begin{tabular}{@{}lll@{}}
\toprule
 & \code{float} & \code{double}\\
\midrule
significant bits $p$ (\code{digits}) & \num{\valDigitsFloat} & \num{\valDigitsDouble}\\
decimal digits (\code{digits10}) & \num{\valDigitsTenFloat} & \num{\valDigitsTenDouble}\\
machine epsilon $\epsmach$ & \num{\valEpsFloat} & \num{\valEpsDouble}\\
largest finite $\fmax$ (\code{max}) & \num{\valMaxFloat} & \num{\valMaxDouble}\\
smallest normal (\code{min}) & \num{\valMinFloat} & \num{\valMinDouble}\\
smallest subnormal (\code{denorm_min}) & \num{\valDenormMinFloat} & \num{\valDenormMinDouble}\\
\bottomrule
\end{tabular}
\end{table}

\section{Machine epsilon, the unit roundoff and the ulp}\label{sec:fp-epsilon}

\emph{Machine epsilon} $\epsmach = 2^{1-p}$ is the distance from 1 to the next larger number. Just above 1 the
exponent is 0, and the last of the $p$ binary digits of the significand is worth $2^{1-p}$. For \code{double} this is
$2^{-52}$. The spacing at a general $x$ is called its \emph{unit in the last place}:
$\ulp(x) = \epsmach \cdot 2^{e}$ for $2^{e} \le \abs{x} < 2^{e+1}$. It is the worth of the last digit, scaled by the
exponent. So the spacing is proportional to $\abs{x}$, within a factor of two. For example, $\ulp(1) =
\num{\valUlpAtOne}$; just below 1, the spacing is \num{\valUlpBelowOne}; and $\ulp(1000) = \num{\valUlpAtOneThousand}$.

IEEE~754 arithmetic rounds the exact result of each operation $+$, $-$, $\times$, $/$ and $\sqrt{\phantom{x}}$ to the
nearest representable number. A tie, a result exactly halfway between two numbers, goes to the one whose last digit
is even. The error is at most half a spacing, so
\begin{equation*}
  \fl(a \circ b) = (a \circ b)(1 + \delta), \qquad \abs{\delta} \le \epsmach / 2,
\end{equation*}
as long as the result is neither too large nor too small for the format. Here $\circ$ stands for any of the
operations, and $\fl$ for the rounded result. The bound $\epsmach / 2$ is called the \emph{unit roundoff}; many texts
call it $u$. This book speaks of $\epsmach$ only, and the library's rounding model uses it too. The reason is that a
value of $f$ that is the result of many operations is rarely accurate to half an ulp. So $\epsmach \abs{f}$ is a more
honest estimate of its error.

\section{The ends of the range}\label{sec:fp-range}

The largest finite number is $\fmax$. A result beyond it overflows: it rounds to infinity. And infinity minus
infinity is NaN, ``not a number''. The library checks every quantity that could overflow. It checks such a quantity
before it is formed, except the final quotient, which it checks after. If a check fails, the library reports
\code{error::non_finite} instead of a result. It never calls \code{fn} at a point that is not finite.

The \emph{normal} numbers are those of the form above, with $1 \le m < 2$. At the other end of the range lies the
smallest of them, \num{\valMinDouble} for \code{double}. Below it, the exponent can fall no further. The numbers continue down to zero with a fixed spacing,
equal to the smallest subnormal, \num{\valDenormMinDouble}. These \emph{subnormal} numbers fill the gap to zero
gradually. But their relative accuracy falls as they shrink, because a fixed spacing is a larger and larger fraction
of a smaller and smaller number. So rounding there is absolute, not relative. The library's rounding model accounts
for that. It takes the rounding error of a nonzero value of $f$ as $\epsmach \abs{f}$, but never less than the
smallest subnormal.

\section{Why $x + h$ is not $x + h$}\label{sec:fp-step}

A finite difference evaluates $f$ at $x + h$, and that sum is rounded like any other. So the function sees
$\fl(x + h)$. This differs from $x + h$ by up to half the spacing at the result. That is $\ulp(x) / 2$, or $\ulp(x)$
when $x + h$ crosses a power of two, where the spacing doubles (the two-sided spacing of \cref{sb:ulp}).

Two consequences follow, and \cref{ch:rounding} meets both. The first is that a step below $\ulp(x) / 2$ does not
move $x$. If neither $x + h$ nor $x - h$ moves, the library reports \code{error::non_finite}. If only one of them
moves, as can happen near a power of two, the outcome depends on the rule. A two-node rule for $f'$ returns a finite
value that is mostly rounding. Every other rule reports \code{non_finite} (\cref{sb:ulp,sec:api-derivative}). Hence
the requirement: a step well above $\ulp(x)$. The second consequence holds for any step. The distance between the two
points of the central difference is not exactly $2h$. That is why the two-node rules divide by the distance actually
evaluated (\cref{sb:actual-distance}).

\section{Cancellation}\label{sec:fp-cancellation}

The subtraction of two nearby numbers of the same sign is exact. Precisely: if $a, b > 0$ and $a / 2 \le b \le 2a$,
then $\fl(a - b) = a - b$. This is Sterbenz's lemma (Theorem 4.3.1), and the same holds for
$a, b < 0$.\cite{sterbenz1974} So the danger is not in the subtraction but in its inputs. Suppose that $a$ and $b$
carry errors $\delta_a$ and $\delta_b$ from earlier rounding. Then the difference carries $\delta_a - \delta_b$, and
by the triangle inequality its relative error is
\begin{equation*}
  \frac{\abs{\delta_a - \delta_b}}{\abs{a - b}} \le \frac{\abs{\delta_a} + \abs{\delta_b}}{\abs{a - b}},
\end{equation*}
which is large when $a$ and $b$ are close. An error that was small compared with $a$ and $b$ is large compared with
their difference. This is \emph{cancellation}. It is the whole reason for the rounding branch of the error model. As
$h$ shrinks, the values $f(x + h)$ and $f(x - h)$ come closer, but their rounding errors stay the same size. And the
division by $2h$ magnifies what is left (\cref{sec:rounding-v}).

\section{Other types}\label{sec:fp-types}

Nothing above is special to binary. Decimal types round to decimal digits, with their own $\epsmach$. Multiprecision
types, which carry more digits than \code{double}, follow the same rules with more digits. An example is
\code{cpp_bin_float_50} of Boost.Multiprecision, with $p = \num{\valDigitsMultiprecision}$ bits and
$\epsmach = \num{\valEpsMultiprecision}$. What the library needs from a type is its arithmetic, its comparisons and a
few members of \code{std::numeric_limits}; \cref{sec:api-real} lists them. The library's error model needs only
$\epsmach$.
